\documentclass[12pt,a4paper]{article}
\usepackage[utf8]{inputenc}
\usepackage[indonesian]{babel}
\usepackage{amsmath}
\usepackage{amsfonts}
\usepackage{amssymb}
\usepackage{geometry}

\geometry{margin=2.5cm}

\title{Menghitung Panjang Busur Kurva}
\author{Ahmad Fakhrul Bawani}
\date{}

\begin{document}

\maketitle

\section*{Soal}

Find the length of the following curve:
\begin{equation*}
  y = \frac{3}{4}x^{\frac{4}{3}} - \frac{3}{8}x^{\frac{2}{3}} + 5, \quad \text{untuk } 1 \le x \le 27
\end{equation*}

\subsection*{1. Rumus Umum}
Rumus untuk menghitung panjang busur (arc length) fungsi dalam variabel $x$ dari $x = a$ sampai $x = b$ adalah:
\begin{equation*}
  L = \int_{a}^{b} \sqrt{1 + \left(\frac{dy}{dx}\right)^2} \, dx
\end{equation*}

\subsection*{2. Menentukan Turunan Fungsi}
Cari turunan pertama $y$ terhadap $x$:
\begin{align*}
  \frac{dy}{dx} &= \frac{3}{4} \cdot \left(\frac{4}{3}x^{\frac{1}{3}}\right) - \frac{3}{8} \cdot \left(\frac{2}{3}x^{-\frac{1}{3}}\right) + 0 \\
  \frac{dy}{dx} &= x^{\frac{1}{3}} - \frac{1}{4}x^{-\frac{1}{3}}
\end{align*}

Selanjutnya, kuadratkan hasil turunan tersebut:
\begin{align*}
  \left(\frac{dy}{dx}\right)^2 &= \left(x^{\frac{1}{3}} - \frac{1}{4}x^{-\frac{1}{3}}\right)^2 \\
  &= \left(x^{\frac{1}{3}}\right)^2 - 2\left(x^{\frac{1}{3}}\right)\left(\frac{1}{4}x^{-\frac{1}{3}}\right) + \left(\frac{1}{4}x^{-\frac{1}{3}}\right)^2 \\
  &= x^{\frac{2}{3}} - \frac{1}{2} + \frac{1}{16}x^{-\frac{2}{3}}
\end{align*}

Tambahkan dengan 1 untuk dimasukkan ke dalam akar:
\begin{align*}
  1 + \left(\frac{dy}{dx}\right)^2 &= 1 + x^{\frac{2}{3}} - \frac{1}{2} + \frac{1}{16}x^{-\frac{2}{3}} \\
  &= x^{\frac{2}{3}} + \frac{1}{2} + \frac{1}{16}x^{-\frac{2}{3}}
\end{align*}

Perhatikan bahwa bentuk di atas kembali membentuk kuadrat sempurna:
\begin{equation*}
  x^{\frac{2}{3}} + \frac{1}{2} + \frac{1}{16}x^{-\frac{2}{3}} = \left(x^{\frac{1}{3}} + \frac{1}{4}x^{-\frac{1}{3}}\right)^2
\end{equation*}

\subsection*{3. Solusi Integrasi}
Substitusikan ke dalam rumus panjang busur dengan batas $a = 1$ dan $b = 27$:
\begin{align*}
  L &= \int_{1}^{27} \sqrt{\left(x^{\frac{1}{3}} + \frac{1}{4}x^{-\frac{1}{3}}\right)^2} \, dx \\
  L &= \int_{1}^{27} \left(x^{\frac{1}{3}} + \frac{1}{4}x^{-\frac{1}{3}}\right) \, dx
\end{align*}

Lakukan integrasi terhadap masing-masing suku:
\begin{align*}
  L &= \left[ \frac{3}{4}x^{\frac{4}{3}} + \frac{1}{4} \cdot \left(\frac{3}{2}x^{\frac{2}{3}}\right) \right]_{1}^{27} \\
  L &= \left[ \frac{3}{4}x^{\frac{4}{3}} + \frac{3}{8}x^{\frac{2}{3}} \right]_{1}^{27}
\end{align*}

Masukkan batas atas ($x = 27$) dan batas bawah ($x = 1$):
\begin{itemize}
  \item Untuk $x = 27$:
    \begin{align*}
      27^{\frac{4}{3}} &= (3^3)^{\frac{4}{3}} = 3^4 = 81 \\
      27^{\frac{2}{3}} &= (3^3)^{\frac{2}{3}} = 3^2 = 9 \\
      \text{Nilai} &= \frac{3}{4}(81) + \frac{3}{8}(9) = \frac{243}{4} + \frac{27}{8} = \frac{486}{8} + \frac{27}{8} = \frac{513}{8}
    \end{align*}
  \item Untuk $x = 1$:
    \begin{align*}
      \text{Nilai} &= \frac{3}{4}(1) + \frac{3}{8}(1) = \frac{3}{4} + \frac{3}{8} = \frac{6}{8} + \frac{3}{8} = \frac{9}{8}
    \end{align*}
\end{itemize}

Hitung selisihnya:
\begin{align*}
  L &= \frac{513}{8} - \frac{9}{8} \\
  L &= \frac{504}{8} \\
  L &= 63
\end{align*}

Jadi, panjang busur kurva tersebut adalah \textbf{$63$ satuan panjang}.

\end{document}